% GENERATED FILE. Edit canonical lessons, then run npm run generate:books.
% canonical-source-hash: fnv1a64:4833eac6204801da
% canonical-lessons: JA-C177-karai, JA-C177-suppai, JA-C177-mazui, JA-C177-kogeru, JA-C177-kusaru, JA-R177-review-172-177

\chapter{Spicy, Sour and Burnt}
\label{ch:ja-spicy-sour-and-burnt}
\hlchaptermodality{177}

\begin{chapteropening}
\textbf{By the end of this chapter:} \emph{I can say how food tastes and when it goes bad.}

Complete the last lesson of chapter 177: I can say how food tastes and when it goes bad.
\end{chapteropening}

\section[\texorpdfstring{\ja{からい} (karai)}{からい (karai)}]{\ja{からい} (karai) — spicy; hot, salty}

\label{lesson:JA-C177-karai}



\begin{quote}
Before the new one: say the Japanese for to steam, then the Japanese for to stir-fry.
\end{quote}

\subsection*{You'll want to know: \ja{からい}}
\textbf{\ja{からい}} — \emph{karai} — \textquotedblleft{}spicy; hot, salty\textquotedblright{}.

\begin{quote}
\textbf{\ja{この} \ja{りょうりは} \ja{からいです。}} — \emph{kono ryōri wa karai desu} — \textquotedblleft{}This dish is spicy.\textquotedblright{}
\end{quote}

\begin{cousinweb}
\textbf{\ja{からい}} (\emph{karai}) is a native Japanese word (wago), inherited rather than borrowed from Chinese; its kanji was attached later as a kun'yomi (JMdict, KANJIDIC2).
\end{cousinweb}

\begin{grammarlens}[title={What you've built}]
One more word for everyday life at A2.
\end{grammarlens}

\subsection*{Your turn}
\begin{itemize}
\raggedright
  \item \emph{Say it:} \emph{karai}
  \item \emph{Say it:} \emph{karai}, once more
  \item \emph{Say it:} \emph{kono ryōri wa karai desu} — This dish is spicy.
  \item \emph{Recall:} say the Japanese for to steam, then the Japanese for to stir-fry, then say \emph{karai} again
\end{itemize}

\begin{tcolorbox}[breakable,colback=teal!4,colframe=teal!35!black,title={Before you move on}]
What does \ja{からい} mean? (\textquotedblleft{}Spicy; hot, salty\textquotedblright{}.) Say it once more.
\end{tcolorbox}
\section[\texorpdfstring{\ja{すっぱい} (suppai)}{すっぱい (suppai)}]{\ja{すっぱい} (suppai) — sour}

\label{lesson:JA-C177-suppai}



\begin{quote}
Before the new one: say the Japanese for to stir-fry, then the Japanese for spicy.
\end{quote}

\subsection*{You'll want to know: \ja{すっぱい}}
\textbf{\ja{すっぱい}} — \emph{suppai} — \textquotedblleft{}sour\textquotedblright{}.

\begin{quote}
\textbf{\ja{この} \ja{みかんは} \ja{すっぱいです。}} — \emph{kono mikan wa suppai desu} — \textquotedblleft{}This mandarin is sour.\textquotedblright{}
\end{quote}

\begin{cousinweb}
\textbf{\ja{すっぱい}} (\emph{suppai}) is a native Japanese word (wago), inherited rather than borrowed from Chinese; its kanji was attached later as a kun'yomi (JMdict, KANJIDIC2).
\end{cousinweb}

\begin{grammarlens}[title={What you've built}]
One more word for everyday life at A2.
\end{grammarlens}

\subsection*{Your turn}
\begin{itemize}
\raggedright
  \item \emph{Say it:} \emph{suppai}
  \item \emph{Say it:} \emph{suppai}, once more
  \item \emph{Say it:} \emph{kono mikan wa suppai desu} — This mandarin is sour.
  \item \emph{Recall:} say the Japanese for to stir-fry, then the Japanese for spicy, then say \emph{suppai} again
\end{itemize}

\begin{tcolorbox}[breakable,colback=teal!4,colframe=teal!35!black,title={Before you move on}]
What does \ja{すっぱい} mean? (\textquotedblleft{}Sour\textquotedblright{}.) Say it once more.
\end{tcolorbox}
\section[\texorpdfstring{\ja{まずい} (mazui)}{まずい (mazui)}]{\ja{まずい} (mazui) — bad-tasting; unpleasant}

\label{lesson:JA-C177-mazui}



\begin{quote}
Before the new one: say the Japanese for spicy, then the Japanese for sour.
\end{quote}

\subsection*{You'll want to know: \ja{まずい}}
\textbf{\ja{まずい}} — \emph{mazui} — \textquotedblleft{}bad-tasting; unpleasant\textquotedblright{}.

\begin{quote}
\textbf{\ja{この} \ja{コーヒー}\ja{は} \ja{まずいです。}} — \emph{kono kōhī wa mazui desu} — \textquotedblleft{}This coffee tastes bad.\textquotedblright{}
\end{quote}

\begin{cousinweb}
\textbf{\ja{まずい}} (\emph{mazui}) is a native Japanese word (wago), inherited rather than borrowed; JMdict marks it as usually written in kana, and its kanji spelling was chosen for its meaning (JMdict, KANJIDIC2).
\end{cousinweb}

\begin{grammarlens}[title={What you've built}]
One more word for everyday life at A2.
\end{grammarlens}

\subsection*{Your turn}
\begin{itemize}
\raggedright
  \item \emph{Say it:} \emph{mazui}
  \item \emph{Say it:} \emph{mazui}, once more
  \item \emph{Say it:} \emph{kono kōhī wa mazui desu} — This coffee tastes bad.
  \item \emph{Recall:} say the Japanese for spicy, then the Japanese for sour, then say \emph{mazui} again
\end{itemize}

\begin{tcolorbox}[breakable,colback=teal!4,colframe=teal!35!black,title={Before you move on}]
What does \ja{まずい} mean? (\textquotedblleft{}Bad-tasting; unpleasant\textquotedblright{}.) Say it once more.
\end{tcolorbox}
\section[\texorpdfstring{\ja{こげる} (kogeru)}{こげる (kogeru)}]{\ja{こげる} (kogeru) — to get burnt; to scorch}

\label{lesson:JA-C177-kogeru}



\begin{quote}
Before the new one: say the Japanese for sour, then the Japanese for bad-tasting.
\end{quote}

\subsection*{You'll want to know: \ja{こげる}}
\textbf{\ja{こげる}} — \emph{kogeru} — \textquotedblleft{}to get burnt; to scorch\textquotedblright{}.

\begin{quote}
\textbf{\ja{さかなが} \ja{よく} \ja{こげます。}} — \emph{sakana ga yoku kogemasu} — \textquotedblleft{}The fish often gets burnt.\textquotedblright{}
\end{quote}

Polite \textbf{\ja{こげます}}, \emph{kogemasu}; an ichidan verb.

\begin{cousinweb}
\textbf{\ja{こげる}} (\emph{kogeru}) is a native Japanese word (wago), inherited rather than borrowed from Chinese; its kanji was attached later as a kun'yomi (JMdict, KANJIDIC2).
\end{cousinweb}

\begin{grammarlens}[title={What you've built}]
One more word for everyday life at A2.
\end{grammarlens}

\subsection*{Your turn}
\begin{itemize}
\raggedright
  \item \emph{Say it:} \emph{kogeru}
  \item \emph{Say it:} \emph{kogeru}, once more
  \item \emph{Say it:} \emph{sakana ga yoku kogemasu} — The fish often gets burnt.
  \item \emph{Recall:} say the Japanese for sour, then the Japanese for bad-tasting, then say \emph{kogeru} again
\end{itemize}

\begin{tcolorbox}[breakable,colback=teal!4,colframe=teal!35!black,title={Before you move on}]
What does \ja{こげる} mean? (\textquotedblleft{}To get burnt; to scorch\textquotedblright{}.) Say it once more.
\end{tcolorbox}
\section[\texorpdfstring{\ja{くさる} (kusaru)}{くさる (kusaru)}]{\ja{くさる} (kusaru) — to go bad; to rot, to spoil}

\label{lesson:JA-C177-kusaru}



\begin{quote}
Before the new one: say the Japanese for bad-tasting, then the Japanese for to get burnt.
\end{quote}

\subsection*{You'll want to know: \ja{くさる}}
\textbf{\ja{くさる}} — \emph{kusaru} — \textquotedblleft{}to go bad; to rot, to spoil\textquotedblright{}.

\begin{quote}
\textbf{\ja{なつは} \ja{たべものが} \ja{よく} \ja{くさります。}} — \emph{natsu wa tabemono ga yoku kusarimasu} — \textquotedblleft{}In summer food often goes bad.\textquotedblright{}
\end{quote}

Polite \textbf{\ja{くさります}}, \emph{kusarimasu}; a godan verb.

\begin{cousinweb}
\textbf{\ja{くさる}} (\emph{kusaru}) is a native Japanese word (wago), inherited rather than borrowed from Chinese; its kanji was attached later as a kun'yomi (JMdict, KANJIDIC2).
\end{cousinweb}

\begin{grammarlens}[title={What you've built}]
One more word for everyday life at A2.
\end{grammarlens}

\subsection*{Your turn}
\begin{itemize}
\raggedright
  \item \emph{Say it:} \emph{kusaru}
  \item \emph{Say it:} \emph{kusaru}, once more
  \item \emph{Say it:} \emph{natsu wa tabemono ga yoku kusarimasu} — In summer food often goes bad.
  \item \emph{Recall:} say the Japanese for bad-tasting, then the Japanese for to get burnt, then say \emph{kusaru} again
\end{itemize}

\begin{tcolorbox}[breakable,colback=teal!4,colframe=teal!35!black,title={Before you move on}]
What does \ja{くさる} mean? (\textquotedblleft{}To go bad; to rot, to spoil\textquotedblright{}.) Say it once more.
\end{tcolorbox}
\section[(dialogue)]{Kitchen and table words, a review}

\label{lesson:JA-R177-review-172-177}



\begin{quote}
Say the Japanese for meal and lunch.
\end{quote}

\subsection*{Your turn}
Say these aloud:

\begin{itemize}
\raggedright
  \item meal, lunch, to get hungry, dinner, to grill
  \item food, beverage, boxed lunch, to be left over, to savour
  \item pot, kitchen knife, to chop, plate, to dish up
  \item cooking, to boil food, to mix, rice bowl, oil
  \item soy sauce, soybean paste, to crack, to steam, to stir-fry
  \item spicy, sour, bad-tasting, to get burnt, to go bad
\end{itemize}

\begin{tcolorbox}[breakable,colback=teal!4,colframe=teal!35!black,title={Before you move on}]
Meal? (\textbf{\ja{しょくじ}}, \emph{shokuji}.) Lunch? (\textbf{\ja{ひるごはん}}, \emph{hirugohan}.) To get hungry? (\textbf{\ja{すく}}, \emph{suku}.) Dinner? (\textbf{\ja{ばんごはん}}, \emph{bangohan}.) To grill? (\textbf{\ja{やく}}, \emph{yaku}.) Food? (\textbf{\ja{たべもの}}, \emph{tabemono}.) Beverage? (\textbf{\ja{のみもの}}, \emph{nomimono}.) Boxed lunch? (\textbf{\ja{おべんとう}}, \emph{obentō}.) To be left over? (\textbf{\ja{のこる}}, \emph{nokoru}.) To savour? (\textbf{\ja{あじわう}}, \emph{ajiwau}.) Pot? (\textbf{\ja{なべ}}, \emph{nabe}.) Kitchen knife? (\textbf{\ja{ほうちょう}}, \emph{hōchō}.) To chop? (\textbf{\ja{きざむ}}, \emph{kizamu}.) Plate? (\textbf{\ja{おさら}}, \emph{osara}.) To dish up? (\textbf{\ja{もる}}, \emph{moru}.) Cooking? (\textbf{\ja{りょうり}}, \emph{ryōri}.) To boil food? (\textbf{\ja{ゆでる}}, \emph{yuderu}.) To mix? (\textbf{\ja{まぜる}}, \emph{mazeru}.) Rice bowl? (\textbf{\ja{ちゃわん}}, \emph{chawan}.) Oil? (\textbf{\ja{あぶら}}, \emph{abura}.) Soy sauce? (\textbf{\ja{しょうゆ}}, \emph{shōyu}.) Soybean paste? (\textbf{\ja{みそ}}, \emph{miso}.) To crack? (\textbf{\ja{わる}}, \emph{waru}.) To steam? (\textbf{\ja{むす}}, \emph{musu}.) To stir-fry? (\textbf{\ja{いためる}}, \emph{itameru}.) Spicy? (\textbf{\ja{からい}}, \emph{karai}.) Sour? (\textbf{\ja{すっぱい}}, \emph{suppai}.) Bad-tasting? (\textbf{\ja{まずい}}, \emph{mazui}.) To get burnt? (\textbf{\ja{こげる}}, \emph{kogeru}.) To go bad? (\textbf{\ja{くさる}}, \emph{kusaru}.)
\end{tcolorbox}
